% Options for packages loaded elsewhere
\PassOptionsToPackage{unicode}{hyperref}
\PassOptionsToPackage{hyphens}{url}
\documentclass[
]{article}
\usepackage{xcolor}
\usepackage[margin=1in]{geometry}
\usepackage{amsmath,amssymb}
\setcounter{secnumdepth}{-\maxdimen} % remove section numbering
\usepackage{iftex}
\ifPDFTeX
  \usepackage[T1]{fontenc}
  \usepackage[utf8]{inputenc}
  \usepackage{textcomp} % provide euro and other symbols
\else % if luatex or xetex
  \usepackage{unicode-math} % this also loads fontspec
  \defaultfontfeatures{Scale=MatchLowercase}
  \defaultfontfeatures[\rmfamily]{Ligatures=TeX,Scale=1}
\fi
\usepackage{lmodern}
\ifPDFTeX\else
  % xetex/luatex font selection
\fi
% Use upquote if available, for straight quotes in verbatim environments
\IfFileExists{upquote.sty}{\usepackage{upquote}}{}
\IfFileExists{microtype.sty}{% use microtype if available
  \usepackage[]{microtype}
  \UseMicrotypeSet[protrusion]{basicmath} % disable protrusion for tt fonts
}{}
\makeatletter
\@ifundefined{KOMAClassName}{% if non-KOMA class
  \IfFileExists{parskip.sty}{%
    \usepackage{parskip}
  }{% else
    \setlength{\parindent}{0pt}
    \setlength{\parskip}{6pt plus 2pt minus 1pt}}
}{% if KOMA class
  \KOMAoptions{parskip=half}}
\makeatother
\usepackage{graphicx}
\makeatletter
\newsavebox\pandoc@box
\newcommand*\pandocbounded[1]{% scales image to fit in text height/width
  \sbox\pandoc@box{#1}%
  \Gscale@div\@tempa{\textheight}{\dimexpr\ht\pandoc@box+\dp\pandoc@box\relax}%
  \Gscale@div\@tempb{\linewidth}{\wd\pandoc@box}%
  \ifdim\@tempb\p@<\@tempa\p@\let\@tempa\@tempb\fi% select the smaller of both
  \ifdim\@tempa\p@<\p@\scalebox{\@tempa}{\usebox\pandoc@box}%
  \else\usebox{\pandoc@box}%
  \fi%
}
% Set default figure placement to htbp
\def\fps@figure{htbp}
\makeatother
\setlength{\emergencystretch}{3em} % prevent overfull lines
\providecommand{\tightlist}{%
  \setlength{\itemsep}{0pt}\setlength{\parskip}{0pt}}
\usepackage{bookmark}
\IfFileExists{xurl.sty}{\usepackage{xurl}}{} % add URL line breaks if available
\urlstyle{same}
\hypersetup{
  pdftitle={Lab 3\_report.rmd},
  pdfauthor={Sarah Wicks},
  hidelinks,
  pdfcreator={LaTeX via pandoc}}

\title{Lab 3\_report.rmd}
\author{Sarah Wicks}
\date{2026-09-14}

\begin{document}
\maketitle

\subsection{Lab 3: Reproducible Research in
R}\label{lab-3-reproducible-research-in-r}

Checkpoint 4: How do we interpret this summary?

The summary table for the ANOVA shows the degrees of freedom, summary of
squares, mean of squares (variance), F value (ratio of between-group to
within-group variance) and the p values associated with each comparison.
None of the 3 p values generated meet the significance threshold (p =
0.05), indicating that the difference between the means is not
statistically significant. In this scenario, the effects of both
population of origin and temperature treatment on the growth rate are
not statistically significant.

Checkpoint 5. What did we just do? Do you think this approach is easily
editable and reproducible?

In order to compare the effect of the temperature treatments and
population of origin on the growth rate, we created a boxplot showing
the growth rate as a function of both temperature treatments and both
populations of origin. To determine whether the effects were
statistically significant, we ran an ANOVA to calculate the difference
between the means of each category. I do think the approach is editable
and reproducible - what's currently lacking in the workflow would be a
Readme to explain the full process.

```\{EditDataframe.R\}

\#\#Import data

temp.df \textless- read.csv(``data/tempExperiment-raw.csv'', header = T)
temp.df

head(temp.df)

print(temp.df)

\#separate temperature treatment from population using strsplit a
\textless- strsplit(as.character(temp.df\(temp), split = '-')
#use unlist and matrix functions to rearranging the output from strsplit
newvar <- matrix(unlist(a), ncol = 2, byrow = TRUE)
head(newvar)
#make a new dataframe with just temp and a new column for pop
df.t2 <- temp.df #set the new dataframe to the original so we can fix it
df.t2\)temp \textless- newvar{[},1{]} df.t2\$pop \textless-
newvar{[},2{]} head(df.t2)

\#use sub to give temp only numeric values and pop only levels pop 1 and
pop 2, watch the space at the end of population
df.t2\(temp <- sub('ten', '10', df.t2\)temp)
df.t2\(temp <- sub('twenty', '20', df.t2\)temp)
df.t2\(pop <- sub('population ', 'pop', df.t2\)pop) str(df.t2) \#Make
sure all objects are correct flavor using factor() or as.numeric
df.t2\(temp <- factor(df.t2\)temp) df.t2\(pop <- factor(df.t2\)pop)
str(df.t2)

\#Export a new edited csv write.csv(df.t2,
`data/tempExperiment\_v2.csv')

\begin{verbatim}



```{Analysis.R}

#start up from last step in EditDataframe.R, which was to export a CSV of our cleaned up data. 
#Clear the environment
rm(list=ls(all=TRUE))
#Import the cleaned data as a new df
df.t <- read.csv('data/tempExperiment_v2.csv')
#Get an overview
str(df.t)

#R thinks temp is a nubmer but really it's a low or high temp category, so change temp to a factor
df.t$temp <- factor(df.t$temp)

##want to analyze how temperature treatment, population of origin, and size affect growth rate


#Plot the data - boxplot of variable growth rate as a function of (~) temp and pop. + tells it to combine temp and pop to create 4 distinct groups (ex. for df$A + df$b then it does A&C, B&C, A&D, B&D)
boxplot(df.t$growthRate~df.t$temp+df.t$pop)

#add axis labels

boxplot(df.t$growthRate~df.t$temp+df.t$pop,
        names = c('10', '20', '10', '20'), #temperature label
        at = c(1,2,4,5), #space out boxes according to population - I don't get this
        ylab = 'Growth rate mm/day',
        xlab = ''
)


#use mtext to add text in the margin - not an axis label b/c you want to space it with the boxplot. idk how these parameters change it

mtext('Pop 1', side = 1, at = 1.5, line = 3)
mtext('Pop 2', side = 1, at = 4.5, line = 3)

#Now save it to PDF in results folder

pdf('results/MyBoxplot.pdf', width = 5, height = 5)
##Use code from above
boxplot(df.t$growthRate~df.t$temp+df.t$pop,
        names = c('10', '20', '10', '20'), #temperature label
        at = c(1,2,4,5), #space out boxes according to population - I don't get this
        ylab = 'Growth rate mm/day',
        xlab = ''
)


#use mtext to add text in the margin - not an axis label b/c you want to space it with the boxplot. idk how these parameters change it

mtext('Pop 1', side = 1, at = 1.5, line = 3)
mtext('Pop 2', side = 1, at = 4.5, line = 3)

dev.off()

##Now find out if effect of pop or effect of temp is significant in analysis. Use ANOVA. 
# The '*' symbol is for multiplication and means we want to evaluate an interaction, in this cas evaluate growth rate as a function of the interaction between temp & pop

m1 <- aov(df.t$growthRate ~ df.t$temp * df.t$pop)
m1.summary <- summary(m1)
m1.summary

##save the summary table as an R object so you can call it later
saveRDS(m1.summary, 'results/m1.summary.rds')
\end{verbatim}

```\{driver.R\} \#\#Run to call the scripts written during analysis

source(`src/EditDataframe.R')

source(`src/Analysis.R')

```

Note that the \texttt{echo\ =\ FALSE} parameter was added to the code
chunk to prevent printing of the R code that generated the plot.

\end{document}
